\chapter{Fase I: la coyuntura de Bogotá en 2013}\label{cap:coyuntura2013}

La primera fase de la investigación se sitúa en 2013, el año en que Bogotá contrató el diseño de una
primera línea de metro subterránea. El propósito es doble: reproducir la decisión con la información
que existía entonces, sin el beneficio de la retrospectiva, y contrastarla después con lo que ocurrió.
Este capítulo describe la coyuntura económica y de movilidad de ese momento y la secuencia de
decisiones que llevó del metro subterráneo al elevado que hoy se construye; el
\cref{cap:evaluacion2013} aplica el modelo con las condiciones de 2013.

\section{La coyuntura económica}

Bogotá llegó a 2013 con una economía en expansión moderada y un mercado laboral en mejora. En 2012 su
producto interno bruto creció 3,7\,\% y representó el 24,7\,\% del PIB nacional \citep{icer2013}; en
2013 el crecimiento de la ciudad fue 0,7 puntos porcentuales inferior al nacional \citep{dane2014pib},
en un año en que el país crecía con fuerza. La tasa de desempleo de la ciudad bajó a 9,0\,\%, medio
punto menos que en 2012 \citep{icer2013}.

Dos rasgos de ese contexto importan para la evaluación de un proyecto de infraestructura de larga vida.
El primero es la expectativa demográfica: las proyecciones vigentes, construidas con el censo de 2005,
asignaban a Bogotá \PobBogotaDosMilTrece{} millones de habitantes en 2013 y un crecimiento cercano a
1,3\,\% anual; para 2018 proyectaban 8.181.047 habitantes. El censo de 2018 encontró 7.181.469, un
millón menos \citep{dane2019censo}. Es decir, en 2013 se planeaba para una ciudad más grande y de
crecimiento más rápido que la que finalmente se observó. El segundo rasgo es la tasa social de
descuento: los proyectos públicos se evaluaban entonces con una tasa real de 12\,\%, sustituida en 2022
por la de 9\,\% \citep{dnp2022}. Una tasa más alta penaliza los proyectos cuyo costo se concentra al
principio y cuyo beneficio llega décadas después, que es el perfil de un metro, y castiga con más
fuerza a la alternativa más costosa.

\section{La movilidad en 2013}

La referencia de movilidad de la época es la Encuesta de Movilidad 2011, la primera que incluyó a los
municipios vecinos \citep{sdm2012}. Según el índice de movilidad de la Secretaría Distrital de
Planeación, construido con esa encuesta, el 41\,\% de los viajes motorizados se hacía en transporte
público colectivo y el 18\,\% en TransMilenio, de modo que el transporte público cubría cerca del 60\,\%
de los viajes motorizados \citep{sdp2013}. Con el reparto de la encuesta (transporte público 30\,\%,
automóvil 10\,\%, taxi 4\,\% y una fracción supuesta de 3\,\% para la motocicleta), la participación
del transporte público entre los modos motorizados que usa este modelo es de \PartTPDosMilTrece{}.

El sistema estaba en transición: el Sistema Integrado de Transporte Público empezaba a reemplazar a
los buses tradicionales. Para la velocidad vial de hora pico de 2013 se supone que guardaba con la de
2025 la misma razón que los promedios diurnos reportados por la Secretaría de Movilidad (unos 24 km/h
hacia 2013 y 23 km/h en 2023), lo que da \VelDosMilTrece{} km/h. Es un supuesto de calibración, no un
dato observado.

\section{De un metro subterráneo a uno elevado}

La \cref{tab:cronologia} resume las decisiones sobre la Primera Línea del Metro de Bogotá (PLMB) desde
2013. Conviene detenerse en tres de ellas.

\begin{table}[ht]
\centering\small
\caption{Cronología de la Primera Línea del Metro de Bogotá.}\label{tab:cronologia}
\begin{tabular}{lp{11.8cm}}
\toprule
Año & Hecho \\
\midrule
2013 & El Distrito contrata con el Consorcio L1 (IDOM, Euroestudios, Cano Jiménez) la ingeniería básica avanzada de una línea subterránea \citep{alcaldia2014metro}. \\
2014 & Entrega del diseño: unos 27 km del Portal de las Américas a la calle 127, 27 estaciones, 100\,\% subterránea, puesta en servicio prevista en 2021 \citep{alcaldia2014metro,fdn2015}. \\
2015 & La FDN y el IDU revisan el diseño con un ejercicio de ingeniería de valor \citep{fdn2015}. \\
2016 & La nueva administración distrital propone una primera línea elevada hasta la calle 72. \\
2017 & Factibilidad del metro elevado; el CONPES 3900 lo declara de importancia estratégica, con un costo de 12,95 billones de pesos de 2017, aporte de la Nación de 9,08 billones (70\,\%) y del Distrito de 3,86 billones \citep{conpes3900,alcaldia2017conpes}. \\
2018 & El censo encuentra un millón de habitantes menos de lo proyectado \citep{dane2019censo}. \\
2019 & Se firma el contrato de concesión 163 de 2019 con Metro Línea 1 S.A.S. \citep{emb2023}. \\
2020 & La pandemia reduce a la mitad la demanda de TransMilenio \citep{silla2026}. \\
2023 & A solicitud del Gobierno nacional se estudia soterrar un tramo de la línea; el estudio preliminar estima mayores costos y un retraso de varios años, y la línea sigue elevada \citep{emb2023,eltiempo2023,narvaez2023}. \\
2028 & Operación comercial prevista \citep{emb2026}. \\
\bottomrule
\end{tabular}
\end{table}

\paragraph{El diseño subterráneo de 2014.} El diseño del Consorcio L1 definió una línea totalmente
subterránea de unos 27 km con 27 estaciones, una velocidad comercial de 35 km/h y una carga máxima
proyectada que pasaba de unos 46.000 pasajeros por hora y sentido en 2021 a 81.000 en 2050; su costo
de ejecución se estimó en 15,1 billones de pesos \citep{fdn2015}. Ese costo, como todo costo de
ingeniería básica, tenía una incertidumbre considerable: no había llegado a la etapa de diseño de
detalle.

\paragraph{El cambio a la tipología elevada.} El argumento central del cambio fue el costo. El
CONPES 3900 respaldó un metro elevado de 12,95 billones de pesos de 2017 hasta la calle 72
\citep{alcaldia2017conpes}, con una capacidad de 72.000 pasajeros por hora y sentido; según el
Departamento Nacional de Planeación, el subterráneo actualizado a 2017 costaría 16,2 billones sin
troncales alimentadoras \citep{semana2017}. Estudios posteriores elevaron la estimación del
subterráneo: un análisis de la Sociedad Colombiana de Ingenieros citado en prensa sitúa el costo del
proyecto de 2014 en 25,3 billones \citep{sci2023}, y la Contraloría destacó que las estaciones
subterráneas habrían costado más de cinco veces las del viaducto \citep{contraloria2022}. Las cifras
no son comparables entre sí sin cuidado (cambian el año de los pesos, el alcance y la longitud), y por
eso este trabajo las reexpresa con una sola fuente de costos unitarios (\cref{sec:costos_metro}).

\paragraph{El intento de soterrar en 2023.} Con la obra contratada y en ejecución, en 2023 se estudió
soterrar parte del trazado. El estudio preliminar de la Empresa Metro y el concesionario evaluó varias
alternativas \citep{emb2023}; según la prensa, el sobrecosto de un tramo subterráneo estaría entre 8 y
17 billones de pesos \citep{eltiempo2023}, y el gerente de la empresa señaló que la primera alternativa
retrasaría la operación seis años, hasta 2034, sin aumentar la capacidad \citep{narvaez2023}. La línea
se mantuvo elevada. El estudio de 2023 es, además, la fuente de los costos unitarios que este trabajo
usa para comparar tipologías.

\section{Lo que se sabía y lo que no}

Desde 2013 se conocían la congestión, el costo aproximado de las dos tipologías y la tasa de descuento.
No se conocían tres hechos que resultarían decisivos: que la población era un millón menor de lo
proyectado, que una pandemia reduciría la demanda de transporte público a la mitad en 2020 y que la
demanda en 2025 seguiría por debajo de la de 2019, y que la primera línea abriría en 2028, siete años
después de lo previsto en el diseño de 2014. La evaluación del \cref{cap:evaluacion2013} separa lo
primero de lo segundo: la evaluación \emph{ex ante} usa solo la información de 2013 y la evaluación
\emph{ex post} repite el ejercicio con la trayectoria observada.
